% from ega4-2-fr.tex
\oldpage[IV-2]{22}
est \emph{injectif} par platitude, donc le morphisme $v$ est \emph{dominant}
(\textbf{I}, 1.2.7), ce qui achève la démonstration.

\begin{corollary}[2.4.10]
\label{IV.2.4.10-fr}
Soient $k$ un corps, $X$, $Y$ des $k$-préschémas~; alors le morphisme
projection $X\times_k Y\to X$ est universellement ouvert. En particulier,
pour toute extension $K$ de $k$, le morphisme projection $X_{(K)}\to X$ est
universellement ouvert.
\end{corollary}

Il suffit d'appliquer (2.4.9) au morphisme structural
$Y\to\operatorname{Spec}(k)$.

\begin{remark}[2.4.11]
\label{IV.2.4.11-fr}
Si $f:X\to Y$ est un morphisme \emph{ouvert}, on sait que, pour \emph{toute} partie
$E$ de $Y$, on a
$f^{-1}(\overline E)=\overline{f^{-1}(E)}$
(Bourbaki, \emph{Top. gén.}, chap.~I, 3\textsuperscript{e} éd., \S~5,
n\textsuperscript{o}~4, prop.~7). Cette remarque s'applique par exemple
lorsque $f$ est un morphisme \emph{plat localement de présentation finie}
(2.4.6), ou un morphisme projection $X\times_k Y\to X$ où $X$, $Y$ sont des
préschémas sur un corps $k$ (2.4.10), et généralise alors (2.3.10).
\end{remark}

\subsection{Permanence des propriétés des Modules par descente fidèlement plate.}

\begin{proposition}[2.5.1]
\label{IV.2.5.1-fr}
Soient $f:X\to Y$ un morphisme, $\mathscr F$ un $\mathcal O_X$-Module
quasi-cohérent, $g:Y'\to Y$ un morphisme fidèlement plat,
$X'=X\times_Y Y'$, $f'=f_{(Y')}:X'\to Y'$,
$\mathscr F'=\mathscr F\otimes_{\mathcal O_Y}\mathcal O_{Y'}$. Pour que
$\mathscr F$ soit plat \textup{(resp.} fidèlement plat) relativement à $f$, il faut et
il suffit que $\mathscr F'$ soit plat \textup{(resp.} fidèlement plat) relativement à
$f'$.
\end{proposition}

Il suffit d'appliquer (2.2.10) en remplaçant $X$, $Y'$, $Z$, $\mathscr F$,
$\mathscr G'$ par $Y'$, $X$, $Y$, $\mathcal O_{Y'}$, $\mathscr F$
respectivement~; on en conclut que pour que $\mathscr F$ soit plat (resp.
fidèlement plat) relativement à $f$, il faut et il suffit que $\mathscr F'$
soit plat (resp. fidèlement plat) relativement à $g\circ f'$. Mais si
$g':X'\to X$ est la projection canonique, $g'$ est fidèlement plat (2.2.13)
et on a $g\circ f'=f\circ g'$~; donc pour que $\mathscr F'$ soit plat (resp.
fidèlement plat) relativement à $f\circ g'$, il faut et il suffit que
$\mathscr F$ le soit relativement à $f$ (2.2.11, (iii)).

\begin{proposition}[2.5.2]
\EGAReaderTarget{ega4.c04.prop.2-5-2}{IV.2.5.2-fr}
Soient $f:X'\to X$ un morphisme fidèlement plat et quasi-compact,
$\mathscr F$ un $\mathcal O_X$-Module quasi-cohérent,
$\mathscr F'=f^*(\mathscr F)$. Considérons, pour un Module quasi-cohérent, la
propriété d'être~:
\begin{enumerate}
\item[\textup{(i)}] de type fini~;
\item[\textup{(ii)}] de présentation finie~;
\item[\textup{(iii)}] localement libre de type fini~;
\item[\textup{(iv)}] localement libre de rang $n$.
\end{enumerate}
Alors, si $\boldsymbol{P}$ désigne une des propriétés précédentes, pour que
$\mathscr F$ possède la propriété $\boldsymbol{P}$, il faut et il suffit que
$\mathscr F'$ la possède.
\end{proposition}

Pour qu'un $\mathcal O_X$-Module quasi-cohérent soit localement libre de type
fini, il faut et il suffit qu'il soit plat sur $X$ et de présentation finie
(Bourbaki, \emph{Alg. comm.}, chap.~\textbf{II}, \S~5, n\textsuperscript{o}~2,
cor.~2 du th.~1, compte tenu de (2.1.2))~; comme $\mathscr F$ est plat sur
$X$ si et seulement si $\mathscr F'$ est plat sur $X'$ en vertu de (2.5.1)
(appliqué en prenant pour $f$ l'identité), on voit que pour prouver la
proposition dans le cas (iii), il suffit de l'avoir fait dans les cas (i) et
(ii)~; il en est de même pour prouver (iv), car
$f^*(\mathcal O_Y^n)=\mathcal O_X^n$, donc si $\mathscr F$ et
$\mathscr F'$ sont localement libres de type fini et $x=f(x')$, le rang de
$\mathscr F$ en $x$ est égal à celui de $\mathscr F'$ en $x'$, et notre
assertion résulte de ce que $f$ est surjectif. Pour traiter les cas (i) et
(ii),
